\section{Índices persistentes}
\subsection{Árbol B+}
\texttt{bplus\_node.py} representa y serializa nodos; \texttt{bplus\_tree.py} implementa descenso, inserción y splits. Los internos contienen punteros a hijos y separadores; las hojas contienen pares clave/RID y enlaces bidireccionales. Se usa el par compuesto clave/RID para ordenar claves secundarias repetidas sin perder registros.

\texttt{\_find\_leaf()} desciende desde la raíz leyendo nodos y conserva el camino de padres. Al desbordar una hoja, \texttt{\_split\_leaf()} separa aproximadamente la mitad y copia al padre la primera clave derecha. Si desborda un interno, se promueve su separador central, se divide y se continúa hacia arriba. Si no hay padre se crea una nueva raíz y se persisten su identificador y altura.

Los rangos localizan la hoja inicial, avanzan por \texttt{next\_leaf\_id} y se detienen al superar el máximo. La capacidad depende del tamaño real de página y del tipo de clave. Se utiliza un máximo común de claves que cabe tanto en hojas como en internos. El borrado elimina el par exacto sin fusionar nodos; los límites de separación siguen siendo válidos. Las hojas vacías permanecen enlazadas.

\subsection{Extendible Hashing}
El índice mantiene buckets con profundidad local y un directorio paginado con profundidad global. \texttt{hash\_directory.py} lee bloques de punteros; el directorio completo nunca se materializa en RAM. Un hash BLAKE2 de 64 bits produce los mismos bits al reabrir el proceso, a diferencia de \texttt{hash(str)} de Python. El prefijo utilizado son los bits bajos según la profundidad global.

Si la profundidad local del bucket lleno coincide con la global, se duplica el directorio físicamente. Luego se divide el bucket, aumenta su profundidad local, se redistribuyen entradas según el nuevo bit y se actualizan las referencias correspondientes. El split también funciona sin duplicación cuando ya existe suficiente profundidad global.

Para claves repetidas o colisiones de hash indistinguibles se usan páginas encadenadas; la profundidad máxima es 20 para limitar el tamaño del directorio. No hay contracción ni merge de buckets. Las páginas vacías de colisiones retiradas en un split no se reciclan; reconstruir el índice recupera ese espacio.

\subsection{Integridad y catálogo}
El catálogo JSON conserva esquemas, organización, tamaño de página y nombres de archivos e índices. Cada tabla incluye un Hash interno para comprobar la unicidad de PRIMARY KEY; el planificador no utiliza ese índice privado como ruta de SELECT. Esto permite demostrar un Full Scan real antes de crear índices públicos, sin volver cuadrática la importación. Sus lecturas y escrituras sí forman parte de las métricas del ejecutor.
